\begin{filecontents*}{aaai24embedded.sty}
\NeedsTeXFormat{LaTeX2e}
\ProvidesPackage{aaai24}[2023/06/26 AAAI 2024 Submission format]
\def\year{2024}
\def\copyright@on{T}
\def\showauthors@on{T}
\def\nocopyright{\gdef\copyright@on{}}
\DeclareOption{submission}{%
  \gdef\copyright@on{}%
  \gdef\showauthors@on{}%
  \long\gdef\pdfinfo #1{\relax}%
}
\ProcessOptions\relax

\setlength\topmargin{-0.25in}
\setlength\oddsidemargin{-0.25in}
\setlength\textheight{9.0in}
\setlength\textwidth{7.0in}
\setlength\columnsep{0.375in}
\newlength\titlebox
\setlength\titlebox{2.25in}
\setlength\headheight{0pt}
\setlength\headsep{0pt}
\thispagestyle{empty}
\pagestyle{empty}
\flushbottom
\twocolumn
\sloppy
\def\addcontentsline#1#2#3{}

\def\copyright@year{\number\year}
\def\copyright@text{Copyright \copyright\space \copyright@year,
Association for the Advancement of Artificial Intelligence (www.aaai.org).
All rights reserved.}
\def\copyrighttext#1{\gdef\copyright@on{T}\gdef\copyright@text{#1}}
\def\copyrightyear#1{\gdef\copyright@on{T}\gdef\copyright@year{#1}}

\def\maketitle{%
  \par%
  \begingroup
    \def\thefootnote{\fnsymbol{footnote}}
    \twocolumn[\@maketitle] \@thanks%
  \endgroup%
  \if T\copyright@on
    \insert\footins{\noindent\footnotesize\copyright@text}%
  \fi%
  \setcounter{footnote}{0}%
  \let\maketitle\relax%
  \let\@maketitle\relax%
  \gdef\@thanks{}%
  \gdef\@author{}%
  \gdef\@title{}%
  \let\thanks\relax%
}
\long\gdef\affiliations #1{ \def \affiliations_{\if T\showauthors@on#1\fi}}

\def\@maketitle{%
  \def\theauthors{\if T\showauthors@on\@author\else Anonymous submission\fi}
  \newcounter{eqfn}\setcounter{eqfn}{0}%
  \newsavebox{\titlearea}
  \sbox{\titlearea}{
    \let\footnote\relax\let\thanks\relax%
    \setcounter{footnote}{0}%
    \def\equalcontrib{%
      \ifnum\value{eqfn}=0%
        \footnote{These authors contributed equally.}%
        \setcounter{eqfn}{\value{footnote}}%
      \else%
        \footnotemark[\value{eqfn}]%
      \fi%
    }%
    \vbox{%
      \hsize\textwidth%
      \linewidth\hsize%
      \vskip 0.625in minus 0.125in%
      \centering%
      {\LARGE\bf \@title \par}%
      \vskip 0.1in plus 0.5fil minus 0.05in%
      {\Large{\textbf{\theauthors\ifhmode\\\fi}}}%
      \vskip .2em plus 0.25fil%
      {\normalsize \affiliations_\ifhmode\\\fi}%
      \vskip .5em plus 2fil%
    }%
  }%
  \newlength\actualheight%
  \settoheight{\actualheight}{\usebox{\titlearea}}%
  \ifdim\actualheight>\titlebox%
    \setlength{\titlebox}{\actualheight}%
  \fi%
  \vbox to \titlebox {%
    \let\footnote\thanks\relax%
    \setcounter{footnote}{0}%
    \def\equalcontrib{%
      \ifnum\value{eqfn}=0%
        \footnote{These authors contributed equally.}%
        \setcounter{eqfn}{\value{footnote}}%
      \else%
        \footnotemark[\value{eqfn}]%
      \fi%
    }%
    \hsize\textwidth%
    \linewidth\hsize%
    \vskip 0.625in minus 0.125in%
    \centering%
    {\LARGE\bf \@title \par}%
    \vskip 0.1in plus 0.5fil minus 0.05in%
    {\Large{\textbf{\theauthors\ifhmode\\\fi}}}%
    \vskip .2em plus 0.25fil%
    {\normalsize \affiliations_\ifhmode\\\fi}%
    \vskip .5em plus 2fil%
  }%
}

\renewenvironment{abstract}{%
  \centerline{\bf Abstract}%
  \vspace{0.5ex}%
  \setlength{\leftmargini}{10pt}%
  \begin{quote}%
    \small%
}{%
  \par%
  \end{quote}%
  \vskip 1ex%
}

\def\pubnote#1{
  \thispagestyle{myheadings}%
  \pagestyle{myheadings}%
  \markboth{#1}{#1}%
  \setlength\headheight{10pt}%
  \setlength\headsep{10pt}%
}

\def\section{\@startsection{section}{1}{\z@}
{-2.0ex plus -0.5ex minus -.2ex}
{3pt plus 2pt minus 1pt}{\Large\bf\centering}}
\def\subsection{\@startsection{subsection}{2}{\z@}
{-2.0ex plus -0.5ex minus -.2ex}
{3pt plus 2pt minus 1pt}{\large\bf\raggedright}}
\def\subsubsection{\@startsection{subparagraph}{3}{\z@}
{-6pt plus -2pt minus -1pt}{-1em}{\normalsize\bf}}
\renewcommand\paragraph{\@startsection{paragraph}{4}{\z@}
{-6pt plus -2pt minus -1pt}{-1em}{\normalsize\bf}}
\setcounter{secnumdepth}{0}

\footnotesep 6.65pt
\skip\footins 9pt plus 4pt minus 2pt
\def\footnoterule{\kern-3pt \hrule width 5pc \kern 2.6pt}
\setcounter{footnote}{0}
\parindent 10pt
\topsep 4pt plus 1pt minus 2pt
\partopsep 1pt plus 0.5pt minus 0.5pt
\itemsep 0.5pt plus 1pt minus 0.5pt
\parsep 2pt plus 1pt minus 0.5pt
\leftmargin 10pt
\leftmargini 13pt
\leftmarginii 10pt
\leftmarginiii 5pt
\leftmarginiv 5pt
\leftmarginv 5pt
\leftmarginvi 5pt
\labelwidth\leftmargini
\advance\labelwidth-\labelsep
\labelsep 5pt

\def\@listi{\leftmargin\leftmargini}
\def\@listii{\leftmargin\leftmarginii
\labelwidth\leftmarginii\advance\labelwidth-\labelsep
\topsep 2pt plus 1pt minus 0.5pt
\parsep 1pt plus 0.5pt minus 0.5pt
\itemsep \parsep}
\def\@listiii{\leftmargin\leftmarginiii
\labelwidth\leftmarginiii\advance\labelwidth-\labelsep
\topsep 1pt plus 0.5pt minus 0.5pt
\parsep \z@
\partopsep 0.5pt plus 0pt minus 0.5pt
\itemsep \topsep}
\def\@listiv{\leftmargin\leftmarginiv
\labelwidth\leftmarginiv\advance\labelwidth-\labelsep}
\def\@listv{\leftmargin\leftmarginv
\labelwidth\leftmarginv\advance\labelwidth-\labelsep}
\def\@listvi{\leftmargin\leftmarginvi
\labelwidth\leftmarginvi\advance\labelwidth-\labelsep}

\abovedisplayskip 7pt plus2pt minus5pt
\belowdisplayskip \abovedisplayskip
\abovedisplayshortskip 0pt plus3pt
\belowdisplayshortskip 4pt plus3pt minus3pt

\def\normalsize{\@setfontsize\normalsize\@xpt{11}}
\def\small{\@setfontsize\small\@ixpt{10}}
\def\footnotesize{\@setfontsize\footnotesize\@ixpt{10}}
\def\scriptsize{\@setfontsize\scriptsize\@viipt{10}}
\def\tiny{\@setfontsize\tiny\@vipt{7}}
\def\large{\@setfontsize\large\@xipt{12}}
\def\Large{\@setfontsize\Large\@xiipt{14}}
\def\LARGE{\@setfontsize\LARGE\@xivpt{16}}
\def\huge{\@setfontsize\huge\@xviipt{20}}
\def\Huge{\@setfontsize\Huge\@xxpt{23}}

\AtBeginDocument{%
  \@ifpackageloaded{natbib}{%
    \let\cite\citep
    \let\shortcite\citeyearpar
    \setcitestyle{aysep={}}
    \setlength\bibhang{0pt}
    \bibliographystyle{aaai24}
  }{}%
  \@ifpackageloaded{hyperref}{%
    \PackageError{aaai}{You must not use hyperref in AAAI papers.}
    {Remove the hyperref package.}
  }{}%
  \@ifpackageloaded{bbm}{%
    \PackageError{aaai}{You must not use bbm in AAAI papers.}
    {Remove the bbm package.}
  }{}%
  \@ifpackageloaded{authblk}{%
    \PackageError{aaai}{Package authblk is forbidden.}{Remove it.}
  }{}%
  \@ifpackageloaded{balance}{%
    \PackageError{aaai}{Package balance is forbidden.}{Remove it.}
  }{}%
  \@ifpackageloaded{CJK}{%
    \PackageError{aaai}{Package CJK is forbidden.}{Remove it.}
  }{}%
  \@ifpackageloaded{flushend}{%
    \PackageError{aaai}{Package flushend is forbidden.}{Remove it.}
  }{}%
  \@ifpackageloaded{fontenc}{%
    \PackageError{aaai}{Package fontenc is forbidden.}{Remove it.}
  }{}%
  \@ifpackageloaded{fullpage}{%
    \PackageError{aaai}{Package fullpage is forbidden.}{Remove it.}
  }{}%
  \@ifpackageloaded{geometry}{%
    \PackageError{aaai}{Package geometry is forbidden.}{Remove it.}
  }{}%
  \@ifpackageloaded{grffile}{%
    \PackageError{aaai}{Package grffile is forbidden.}{Remove it.}
  }{}%
  \@ifpackageloaded{navigator}{%
    \PackageError{aaai}{Package navigator is forbidden.}{Remove it.}
  }{}%
  \@ifpackageloaded{savetrees}{%
    \PackageError{aaai}{Package savetrees is forbidden.}{Remove it.}
  }{}%
  \@ifpackageloaded{setspace}{%
    \PackageError{aaai}{Package setspace is forbidden.}{Remove it.}
  }{}%
  \@ifpackageloaded{stfloats}{%
    \PackageError{aaai}{Package stfloats is forbidden.}{Remove it.}
  }{}%
  \@ifpackageloaded{tabu}{%
    \PackageError{aaai}{Package tabu is forbidden.}{Remove it.}
  }{}%
  \@ifpackageloaded{titlesec}{%
    \PackageError{aaai}{Package titlesec is forbidden.}{Remove it.}
  }{}%
  \@ifpackageloaded{tocbibind}{%
    \PackageError{aaai}{Package tocbibind is forbidden.}{Remove it.}
  }{}%
  \@ifpackageloaded{ulem}{%
    \PackageError{aaai}{Package ulem is forbidden.}{Remove it.}
  }{}%
  \@ifpackageloaded{wrapfig}{%
    \PackageError{aaai}{Package wrapfig is forbidden.}{Remove it.}
  }{}%
}
\let\endthebibliography=\endlist
\end{filecontents*}

\documentclass[letterpaper]{article}
\usepackage{aaai24embedded}

\ifdefined\XeTeXversion
  \def\encodingdefault{OT1}
\fi

\usepackage{times,helvet,courier}
\usepackage[hyphens]{url}
\usepackage{graphicx,natbib,caption}
\usepackage{amsmath,amssymb,booktabs,tikz}
\usetikzlibrary{arrows.meta,positioning}

\urlstyle{rm}
\def\UrlFont{\rm}
\nocopyright

\title{Paired Relighting and Routed Consistency\\
for Industrial Anomaly Detection}

\author{}

\affiliations{
Nanyang Technological University\\
AI6103 Deep Learning and Applications\\[0.5em]
\begin{tabular}{@{}p{6.6cm}ll@{}}
Name & Matriculation number & NTU email\\
Maiti Shrijeet & G2609039F & shrijeet001@e.ntu.edu.sg\\
Anand Aryaman & G2607658L & aryaman006@e.ntu.edu.sg\\
Boban Athul Krishna & G2609955K & ATHULKRI001@e.ntu.edu.sg\\
Gantasala Janardhanarao Harinath Umesh
& G2602604G & UMESH001@e.ntu.edu.sg\\
Harsh Saand & G2503201D & harsh011@e.ntu.edu.sg
\end{tabular}
}

\begin{document}
\maketitle

\section{Introduction}

Lighting changes can cause an inspection model to flag an
undamaged product as defective. We investigate whether training
on aligned original and relit images reduces these false alarms
while preserving detection of real defects. We measure both
false-positive rate and defect recall because a reduction in
score variation alone does not establish better detection.

We build on SimpleNet \citep{liu2023simplenet}, which uses a
frozen visual backbone, a trainable feature adapter and a
discriminator trained with synthetic feature anomalies. We keep
the backbone fixed and apply the additional training objective
to its adapter and discriminator. We first compare identity
pairs, relighting alone and relighting with patch-score
consistency. The original proposal focused on this paired
consistency objective. The later extension adds local and
global synthetic anomalies and investigates whether restricting
the consistency gradient to the adapter changes the trade-off
between false alarms and missed defects.

Our contribution is the paired training implementation and a
consistency loss whose gradient updates only the adapter,
combined with local and global synthetic anomalies. The
underlying augmentation, consistency and synthesis techniques
are established. We compare against augmentation-only,
synthesis-only and ordinary-consistency controls to determine
whether the integration helps and which component accounts for
any change.


\section{Background and Baseline}

Industrial anomaly detection identifies defective products using
normal training images, reducing the need for labelled defect
examples. SimpleNet \citep{liu2023simplenet} extracts patch features
with a pretrained visual backbone, adapts them, and trains a
discriminator to distinguish normal features from synthetic
anomalies.

Our baseline uses a frozen WideResNet-50-2 backbone and features
from layers 2 and 3. The concatenated features are mapped to
1536-dimensional patch representations by a trainable adapter.
A discriminator with hidden dimension 1024 is trained using a
hinge objective. Synthetic feature anomalies are generated by
adding Gaussian noise with standard deviation 0.015. At inference,
the negative discriminator output gives the patch anomaly score,
and the image score is the maximum over patches.

For the paired study, we use a fixed-final checkpoint and raw
prediction scores. We do not select checkpoints using test AUROC
or rescale scores using test-set statistics. The three baseline
conditions are B0 (identity pair), B1 (relit second view without
consistency), and B2 (relit pair with patch-score consistency).
They are intended to separate the effect of relighting from the
additional consistency objective.

The later local/global and routed-consistency work is a separate
follow-up, not part of the baseline implementation. PR \#37 adds
loss functions for local/global synthetic anomalies and
adapter-only score consistency, but does not include an image
synthesis pipeline or a training command. The PR \#39 launcher
loads its trainer and related modules from a separately pinned
execution archive. The inspected Git changes do not establish
that this trainer calls the functions introduced in PR \#37.
We therefore distinguish the paired baseline from the follow-up
implementation and do not treat matching arm names as proof
that the same loss code was used.

\section{Experimental Setup}

The baseline implementation is evaluated on MVTec AD and
MVTec AD 2. MVTec AD provides an implementation check across
15 object and texture categories. Images are resized to 329 pixels
and centre-cropped to 288 pixels. We use batch size 8, seeds
0, 1, and 2, and a fixed-final schedule of 40 meta-epochs with
four discriminator passes. Normal training images are split
using an 80/20 training and validation holdout. Identical RGB
images are grouped to avoid duplicates crossing the split.

The paired lighting study uses the eight public MVTec AD 2
categories: can, fabric, fruit\_jelly, rice, sheet\_metal, vial,
wallplugs, and walnuts. The committed manifests contain 2,528
normal training images and 302 normal validation images. The
public test set contains 1,084 images: 379 normal and 705
defective. Images are resized to $256\times256$ using full-frame
bilinear interpolation, without a centre crop. The baseline
study uses batch size 8 and seeds 0, 1, and 2. Relighting samples
brightness and gamma independently from $[0.8,1.2]$.

For the paired comparison, B0 uses an identity pair, B1 uses
relighting without consistency, and B2 adds patch-score
consistency. The earlier B2 sweep considered
$\lambda\in\{0.01,0.1,1.0\}$; these settings are distinct from
the later local/global follow-up. The primary AD 2 baseline
runs used Python 3.12, PyTorch 2.6.0 with CUDA 12.4, and an
NVIDIA A40 GPU.

The follow-up launchers in PR \#39 describe seven AD 2 arms
(B0, B1, B2, L0, L1, L2, and R0) across three seeds and eight
categories. However, the trainer, optimizer configuration,
and loss implementation are loaded from an external archive,
not supplied in the PR. The launcher records the archive hash
and guards the staged data and completed checkpoints, but the
Git checkout alone is insufficient to reproduce training.
The inspected PRs also do not connect the archive's trainer
to the loss functions in PR \#37. Accordingly, details that
depend on the external trainer, including the exact follow-up
loss weights and their implementation, are not assumed here
unless they are explicitly available in the released source.

The blade follow-up uses a second external execution archive.
Its launcher separates fit, calibration, and development data
from optional test packets and accepts a case for packing only
when it reaches 5,120 updates. The documentation states that
a fresh full-budget run through this wrapper had not been
completed; historical results used the original workers.
The five additional blade arms (L3, R1, R2, P0, and P1) are
not defined by the loss code in PR \#37. Their behaviour
therefore requires confirmation from the archived trainer
rather than inference from their names.

\section{Evaluation Protocol}

For the paired baseline, one image-level threshold is calibrated
per model, category, and seed using regular-light normal
validation scores. It is the higher empirical 95th percentile.
An image is classified as anomalous only when its score is
strictly greater than this threshold. The same threshold is
applied to the other lighting conditions. Test labels are not
used to select the checkpoint, threshold, or consistency weight.

We report image-level false-positive rate (FPR), defect recall,
and image AUROC. FPR is the proportion of normal images
classified as anomalous; recall is the proportion of defective
images correctly detected. AUROC measures the ranking of
normal and defective image scores. Category-macro results give
each category equal weight. Pixel AUROC and AU-PRO are
localisation metrics for the paired study; AU-PRO is integrated
up to a false-positive rate of 0.05.

For MVTec AD 2, regular-light results are reported separately
from the non-regular group, which includes underexposure,
overexposure, and additional-light conditions identified by
the \texttt{shift\_*} suffixes. The non-regular aggregate combines
different lighting changes and should not be interpreted as
one uniform unseen condition. The paired-study recall criterion
requires the macro recall drop relative to the relevant reference
arm to be no greater than five percentage points. A reduction
in false alarms alone is not considered sufficient if recall
falls beyond this limit.

The follow-up launchers retain the same broad calibration rule:
the higher empirical 0.95 quantile of normal calibration scores,
with strict score exceedance. The AD 2 launcher scores an image
using the maximum of the negative discriminator output over
patches. Its runtime imports the evaluation modules from the
external archive, so the checkout's evaluation code is not
necessarily the code used to produce the archived runs.
The blade launcher similarly checks calibration replay and
uses a shared threshold across the blade domains.

The results described in PR \#40 are aggregate results from an
external release. Its verifier checks file hashes, row keys,
counts, rates derived from supplied confusion counts, and
agreement between CSV aggregates and summary JSON. It does
not replay model inference, recompute AUROC from raw scores,
or recompute the published confidence intervals. The release
archive was not available during the PR inspection, so the
verifier and reported intervals were not independently run.
The reported AD 2 R0 result reduces changed-light FPR relative
to B0, but also reduces changed-light recall by 27.30 percentage
points. It therefore fails the stated five-point recall
preservation criterion. The blade results show favourable
reported mean values, but seed-level failures and an interval
including zero remain unresolved. These follow-up results are
thus treated as reported, not independently verified evidence.
No pixel AUROC or AU-PRO results are supplied for these new
runs.

\section{Method}

For the extension below, $E$ denotes the frozen feature
extractor, $A_\theta$ the adapter and $D_\phi$ the
discriminator, whose output represents normality. The baseline
architecture and scoring procedure are described in the
baseline section of the full report.

\subsection{Aligned Relighting}

After shared full-frame spatial preprocessing, RGB pixels
$x\in[0,1]$ form the pair $(x,T(x))$, with
\begin{equation}
T(x)=\operatorname{clip}(b x^\gamma,0,1),
\qquad
b,\gamma\sim\mathcal U(0.8,1.2)
\end{equation}
independently.

Relighting occurs before ImageNet normalization. Both views
have the same spatial coordinates, allowing patch-by-patch
comparison. We record clipping because saturation can erase
information. This simple photometric family cannot simulate
every reflection, shadow or physical light-source change.

\subsection{Initial Paired Objective}

For view $v\in\{x,T(x)\}$, let
$z_v=A_\theta(E(v))$ and
$\widetilde z_v=z_v+\epsilon_v$, with independent Gaussian
noise $\epsilon_v\sim\mathcal N(0,0.015^2I)$. With hinge
margin $m=0.5$, the paired discrimination objective is
\begin{align}
L_{\rm hinge}
=\tfrac12\sum_v\big[
&\operatorname{mean}_u(m-D_\phi(z_{v,u}))_+
\notag\\
&+\operatorname{mean}_u
(m+D_\phi(\widetilde z_{v,u}))_+
\big].
\end{align}

Here $(a)_+=\max(a,0)$. Normal-score consistency compares raw
patch outputs:
\begin{equation}
L_{\rm pair}
=
\operatorname{mean}_u
\big[
D_\phi(z_{x,u})-D_\phi(z_{T(x),u})
\big]^2.
\end{equation}

The initial extension uses
$L_{\rm hinge}+\lambda L_{\rm pair}$.
Both adapter and discriminator receive consistency gradients.
One joint discriminator batch contains normal original,
normal relit, synthetic original and synthetic relit features,
in that order. The four groups share one discriminator forward
pass, so batch normalization uses their combined statistics
rather than separate statistics for each view.

\subsection{Local and Global Synthetic Anomalies}

The follow-up retains the same backbone and adapts local and
global synthesis from GLASS \citep{chen2024glass}. Global
perturbations begin with Gaussian offsets. For 20 steps, a
normalized feature gradient ascends the fake-class binary
cross-entropy with step size $0.001$. Every five steps, offsets
are projected to an annulus with radii $0.5$ and $1.0$. Only
the synthetic feature offsets are optimized during mining;
model parameters and normalization buffers are not updated by
that inner operation.

Local synthesis blends a texture donor into a normal image
inside a Perlin-noise region. The implementation combines two
masks by intersection, union or selection; donor rotation,
optional flipping and photometric changes increase appearance
variation. The normal-image blending weight is drawn from
$\mathcal N(0.5,0.1^2)$ and clipped to $[0.2,0.8]$. A fixed
subset of 256 DTD texture donors supplies the inputs
\citep{cimpoi2014dtd}. Four deterministic local variants per
image cycle over source passes.

Normal and local images receive identical relighting draws.
A pixel receives a synthetic-defect label if any channel
differs from the corresponding normal image by more than
$1/255$ after relighting. A feature-grid location is labelled
anomalous if its $8\times8$ pixel block contains at least one
such pixel, producing a $32\times32$ label grid. These are
synthetic labels, not real-defect annotations.

Using anomaly logits $a=-D_\phi(z)$, normal targets are zero
and synthetic global targets are one. Define $L_{\rm global}$
as the sum of the normal and synthetic mean binary
cross-entropies, averaged over the two views. For local labels
$y_i$, $p_i=\sigma(a_i)$ and
$p_{t,i}=y_ip_i+(1-y_i)(1-p_i)$. Local focal supervision is
\begin{equation}
L_{\rm local}
=
\operatorname{mean}_{i\in H}
(1-p_{t,i})^2\operatorname{BCE}(a_i,y_i),
\end{equation}
where $H$ retains squared probability errors at or above their
empirical median, including ties. The ordinary-consistency
control L2 uses
\begin{equation}
L=L_{\rm global}+L_{\rm local}+0.1L_{\rm pair}.
\end{equation}

The global discrimination forward updates the discriminator's
batch-normalization statistics once. Mining and auxiliary
forwards use running statistics; module modes are restored
afterwards.

\subsection{Routing the Consistency Gradient}

In R0, consistency is evaluated through the current
discriminator with detached parameters and cloned buffers,
denoted $D_{\bar\phi}$:
\begin{equation}
L_{\rm route}
=
\operatorname{mean}_u
\big[
D_{\bar\phi}(z_{x,u})-
D_{\bar\phi}(z_{T(x),u})
\big]^2.
\end{equation}

R0 replaces $L_{\rm pair}$ in the preceding objective with
$L_{\rm route}$. The derivative through the discriminator
with respect to its input remains active. Thus
$\nabla_\theta L_{\rm route}$ can change the adapter while
$\nabla_\phi L_{\rm route}=0$. The classification losses
still update both modules. This prevents the consistency term
from directly changing the discriminator parameters. Changes
to the shared adapter can still affect defect scores, so
recall must be checked experimentally.

\begin{figure}[t]
\centering
\begin{tikzpicture}[
box/.style={
draw,
align=center,
font=\small,
text width=29mm,
minimum height=8mm
},
arr/.style={-{Latex[length=1.5mm]}}
]

\node[box] (data)
{normal image and\\aligned relit image};

\node[box,below=4mm of data] (feat)
{frozen features $E$\\trainable adapter $A_\theta$};

\node[box,below left=6mm and -12mm of feat] (disc)
{live $D_\phi$\\classification losses};

\node[box,below right=6mm and -12mm of feat] (aux)
{detached $D_{\bar\phi}$\\consistency loss};

\draw[arr] (data)--(feat);
\draw[arr] (feat)--(disc);
\draw[arr] (feat)--(aux);

\node[font=\small,align=center,below=3mm of disc]
{updates $A_\theta,D_\phi$};

\node[font=\small,align=center,below=3mm of aux]
{updates $A_\theta$ only};

\end{tikzpicture}

\caption{
R0 changes the gradient route of the auxiliary consistency
term. Both branches use the current discriminator values.
Local and global synthetic inputs also feed the classification
losses; they are omitted here to emphasize routing.
}
\label{fig:route}
\end{figure}

\subsection{Matched Controls}

\begin{table}[t]
\centering
\small
\begin{tabular}{llll}
\toprule
Arm & Synthesis/loss & Pair & Consistency\\
\midrule
B0 & Gaussian/hinge & Identity & None\\
B1 & Gaussian/hinge & Relit & None\\
B2 & Gaussian/hinge & Relit & Full, $0.1$\\
L0 & Local+global & Identity & None\\
L1 & Local+global & Relit & None\\
L2 & Local+global & Relit & Full, $0.1$\\
R0 & Local+global & Relit & Adapter, $0.1$\\
\bottomrule
\end{tabular}
\caption{
The pending seven-arm comparison. B2 here is fixed at $0.1$;
the earlier three-weight B2 sweep is a separate completed study.
}
\label{tab:arms}
\end{table}

B1 versus B0 tests augmentation. L1 versus L0 tests
augmentation with local and global synthesis. L1 versus B1
changes synthesis and the discrimination objective, so it does
not isolate routing. R0 versus L1 tests adding routed
consistency; R0 versus L2 tests routing against full classifier
gradients under otherwise matched synthesis. R0 versus B0
measures the combined recipe, not the contribution of routing
alone.

\section{Inspection Demo}

The existing SimpleNet inspection interface displays model
scores, anomaly maps and a lighting sweep for a supplied image.
It uses each model's normal-validation threshold and a fixed
map scale across the sweep. A threshold-relative score is not
a calibrated probability. For an upload without a known label,
a flag cannot be counted as a false alarm.

The recorded blade gallery and the interactive interface are
separate demonstrations. Recorded predictions must be
identified as such. These demonstrations do not establish
that the routed R0 model has been integrated into live
inference, or that the method improves detection across all
categories.

\section{Method Limitations}

Synthetic brightness/gamma changes do not cover all physical
lighting effects, and synthetic anomaly supports do not
establish real-defect sensitivity. Normal validation cannot
certify recall. Auxiliary routing prevents a direct
discriminator consistency gradient, but not indirect changes
through the shared adapter.

\begin{thebibliography}{}

\bibitem[Chen et~al.(2024)]{chen2024glass}
Chen, Q.; Luo, H.; Lv, C.; and Zhang, Z. 2024.
A Unified Anomaly Synthesis Strategy with Gradient Ascent
for Industrial Anomaly Detection and Localization.
In \textit{European Conference on Computer Vision}.

\bibitem[Cimpoi et~al.(2014)]{cimpoi2014dtd}
Cimpoi, M.; Maji, S.; Kokkinos, I.; Mohamed, S.;
and Vedaldi, A. 2014.
Describing Textures in the Wild.
In \textit{Proceedings of the IEEE Conference on Computer
Vision and Pattern Recognition}.

\bibitem[Liu et~al.(2023)]{liu2023simplenet}
Liu, Z.; Zhou, Y.; Xu, Y.; and Wang, Z. 2023.
SimpleNet: A Simple Network for Image Anomaly Detection
and Localization.
In \textit{Proceedings of the IEEE/CVF Conference on
Computer Vision and Pattern Recognition}, 20402--20411.

\end{thebibliography}

\end{document}